% 04 Apr: AI clean

\section{Concepts of Church Fellowship and the Union in 1890}

Editorial article in “Folkebladet” for September 20, 1893.

We shall dwell here upon those misunderstandings—intentional or unintentional—of the concept “church fellowship,” which among some within the United Church have been set in circulation and threaten to gain accepted standing.

A couple of statements in recent times serve to cast a sharp light upon these congregation-endangering notions. In an issue of “Kirkebladet” there was recently—indeed by stealth—smuggled in the remark that Augsburg was a “private” opposition school.

In “Samfundet” No. 2 it is said of Prof. Oftedal—and it must be applied also to Prof. Sverdrup—that he “until quite recently served as a theological professor.”

Such a mode of expression—apart from the childish insult intended in it—rests upon the view that when Augsburg is no longer the seminary of the United Church, then it is a private institution, and that when Professors Sverdrup (and Oftedal) have resigned as theological professors in the United Church, then they are no longer theological professors at all, however much they continue at Augsburg, as they have now been doing for nearly twenty and seventeen years respectively, before anyone had yet heard mention of the United Church.

Such naïve statements and notions spring again not only from that notion of the significance of the United Church, exaggerated beyond all reason, which the majority’s pastors have set in circulation, but essentially, and in the last instance, from an altogether Catholicizing and anti-congregational concept of church fellowship, which the pastors of the old or synodical tendency with all diligence cultivate and exploit—when it can serve their advantage and strengthen their power over the congregation.

For they regard the “association” as a higher form of the congregation, and especially their own fellowship as the real manifestation of the one Church, make it identical with “the Church,” and regard everything outside it as sectarianism.

But is the fellowship such a most holy thing, that the congregation in it as it were disappears, and is handed over into the control of the association?

No, by no means.

Not even the practical carrying out of this by deciding, at the command of a few men, upon the erection of a new and thereby necessarily oppositional seminary, without asking a single congregation for counsel, does not change the matter in the least.

A church fellowship under our free conditions is nothing other than a free mutual agreement between a greater or lesser number of congregations—the number makes no difference to the matter—to work together in certain common undertakings, as for example a seminary, mission, the publishing of religious books, etc., where either an individual congregation is unable to do so on its own, or where the union of common forces will bring each congregation a proportionately greater benefit. Such a union of congregations for mutual support and help is therefore not any larger congregation or a higher form of the congregation; for the Church’s basic unit is and remains, according to God’s Word, the congregation alone, and no church fellowship, not even the Catholic Church, has, without vigorous contradiction from others, been able to arrogate to itself the right to be regarded as “the Church.”

Such a Catholicizing concept of a church fellowship therefore cannot be continued or put into practice without the greatest damage to the congregation’s independence and Christian liberty, and will, since it has only the purpose of advancing the corporate interests of the clergy and gradually bringing forth the office of bishop, unavoidably press the individual congregation and its spiritual need more and more aside, and let our Norwegian Lutheran Church gradually sink back into the Norwegian Synod’s state-churchliness and dead orthodoxy.

Since therefore the joining together of congregations in a fellowship is essentially a practical arrangement for pooling spiritual resources and saving means, it is likewise not such a most holy institution, into which it is easy to come, but which a congregation cannot leave without making itself guilty of schism, sundering, and sin. If a congregation by free decision enters into a fellowship in the hope of finding there spiritual help and mutual edification, and it finds itself disappointed in this expectation, then it is its incontestable right, which no clergy under papal overreach can obstruct, to withdraw from such a fellowship-connection whenever it judges this to be expedient. And although every congregation both ought and will use caution in the exercise of this right and not in all things look only to its own advantage, but also to the advantage of its fellow congregations, still it becomes a downright step toward Catholicism when a clergy, bound together by artificial means, in any way seeks to arrogate to itself authority to hinder a free Lutheran congregation in this inviolable right to join itself to whatever Lutheran fellowship-connection it finds good.

When now the majority in the United Church, through its leaders, claims such extraordinary importance for its church body, that a seminary incorporated for nearly twenty-five years at once becomes a “private” opposition school as soon as the said leaders, in breach of faith and honor, break their—moreover only half-committed—connection with it, or that its professors cease to be professors when they no longer desire to stand in connection with such leaders, then this is all the more absurd and ludicrous in this case, since the congregations of the United Church have not yet formed any completed fellowship-connection.

In 1890 it made the first essential step toward a full union by choosing a common governing body and adopting a common seminary and common mission work for the congregations and fellowships participating in the work of union; but since other significant conditions were attached to this union, and these have not been fulfilled, and since one of the stipulations, namely concerning a “common seminary,” has even by the said leaders of the majority directly and without right been broken without consultation with the congregations, then it ought to be evident to everyone who has not become crazed by the majority that no union in the full sense of the word has yet come into being, and that every congregation is free, and every one of the three contracting parties as well, on account of the majority’s breach of contract, to withdraw from the work of union at this stage without either causing division or withdrawing. For if it is the congregation’s indisputable right to withdraw even from older fellowships, when its welfare requires it, then obviously there can be no question either of schism or of withdrawal where not even any full fellowship-connection has been brought into being. In that case there can only be question of drawing back until a more satisfactory form of union might possibly become the subject of negotiation.

It is therefore clearly high time that the majority leaders tone down their lofty language about the extraordinary importance of that fellowship which the work of union aimed to bring about. The United Church is not identical with the United States, nor with the Church itself. It has only a single, fully corporate legal existence, and that not by virtue of its organization under the law of the congregation or of the Body of Christ, but simply by virtue of its incorporation under the civil law, and even that of a doubtful sort. It thus has the same significance as every other incorporated institution and thereby has—formally—received the right, of which it also has made ample use, namely to initiate legal proceedings. This right any three persons or three congregations whatsoever can acquire by being incorporated under the laws of the state, and is therefore not something over which to exalt oneself so extravagantly.

A measure of sobriety in the concept of fellowship, and modesty in self-assessment, one may demand even of old church fellowships and schools with a fairly fruitful past; the lack of these becomes then not only unchristian, but downright ludicrous in a fellowship which after two years of cooperation has still not attained full union, or in a seminary scarcely more than a few days old.

Sven Oftedal
